% !TEX program = xelatex
\documentclass[11pt,a4paper]{article}

\usepackage[a4paper,margin=2.6cm,top=2.8cm,bottom=2.8cm]{geometry}
\usepackage{fontspec}
\usepackage{microtype}
\usepackage{amsmath,amssymb}
\usepackage{graphicx}
\usepackage{xcolor}
\usepackage{titlesec}
\usepackage{fancyhdr}
\usepackage{hyperref}
\usepackage{bookmark}
\usepackage{etoolbox}
\usepackage{enumitem}

\defaultfontfeatures{Ligatures=TeX}
\setlength{\parindent}{0pt}
\setlength{\parskip}{0.75em}
\setlist{nosep}

\definecolor{IberoTeal}{HTML}{0D9488}
\definecolor{IberoBlue}{HTML}{2563EB}
\definecolor{IberoSlate}{HTML}{334155}
\definecolor{IberoMuted}{HTML}{64748B}
\definecolor{IberoRule}{HTML}{CBD5E1}

\hypersetup{
    colorlinks=true,
    hypertexnames=false,
    linkcolor=IberoBlue,
    urlcolor=IberoTeal,
    pdftitle={IberoSing 2026 Book of Abstracts},
    pdfauthor={IberoSing International Workshop 2026}
}

\titleformat{\section}
    {\Large\bfseries\sffamily\color{IberoTeal}}
    {}
    {0pt}
    {}

\titleformat{\subsection}
    {\large\bfseries\sffamily\color{IberoSlate}}
    {}
    {0pt}
    {}

\pagestyle{fancy}
\setlength{\headheight}{22pt}
\fancyhf{}
\lhead{\sffamily\small IberoSing 2026}
\rhead{\sffamily\small Book of Abstracts}
\cfoot{\sffamily\small\thepage}
\renewcommand{\headrulewidth}{0.4pt}
\renewcommand{\headrule}{\hbox to\headwidth{\color{IberoRule}\leaders\hrule height \headrulewidth\hfill}}

\newcommand{\EntryHeading}[3]{%
    \phantomsection
    \addcontentsline{toc}{section}{#1: #2}
    {\sffamily\Large\bfseries\color{IberoTeal}#2\par}
    \vspace{0.35em}
    {\large\bfseries #1\par}
    \ifstrempty{#3}{}{%
        {\itshape\color{IberoMuted}#3\par}
    }
    \vspace{0.9em}
    {\sffamily\small\bfseries\color{IberoBlue}Abstract\par}
    \vspace{0.2em}
}

\newcommand{\AbstractEntry}[4]{%
    \clearpage
    \EntryHeading{#1}{#2}{#3}
    #4
}

\begin{document}

\begin{titlepage}
    \thispagestyle{empty}
    \centering
    \vspace*{0.5cm}
    \includegraphics[width=\textwidth]{assets/banner.png}

    \vfill

    {\sffamily\bfseries\fontsize{30}{36}\selectfont Book of Abstracts\par}
    \vspace{0.4cm}
    {\sffamily\Large\color{IberoTeal}Zaragoza Week\par}
    \vspace{0.9cm}
    {\Large IberoSing International Workshop 2026\par}
    \vspace{0.35cm}
    {\large June 29--July 3, 2026\par}
    {\large Universidad de Zaragoza, Spain\par}

    \vfill

    {\sffamily\small\color{IberoMuted}
    Fifth edition of the IberoSing International Workshop on Singularity Theory\par}
\end{titlepage}

\pagenumbering{roman}
\tableofcontents
\clearpage
\pagenumbering{arabic}

\section*{Zaragoza Abstracts}
\addcontentsline{toc}{section}{Zaragoza Abstracts}

\AbstractEntry
{Elvira Pérez Callejo}
{Explicit cones for surfaces from pencils at infinity}
{Universidad de Valladolid (IMUVA)}
{
We study rational surfaces obtained by blowing up the projective plane at configurations
$\mathcal{C}=\mathcal{B}\cup \mathcal{D}$, where $\mathcal{B}$ is the set of proper and
infinitely near base points of the pencil associated to a curve with one place at infinity,
and $\mathcal{D}$ is a finite set of free infinitely near points lying on strict transforms
of curves of the pencil. On the resulting surface $X$, the cone of curves
$\operatorname{NE}(X)$ and the nef cone $\operatorname{Nef}(X)$ are polyhedral, and we give
an explicit description of them. Moreover, when the pencil arises from an AMS-type curve and
$\mathcal{D}$ contains at most two free points on each considered curve, the Cox ring
$\operatorname{Cox}(X)$ is finitely generated.

This is joint work with Carlos Galindo, Francisco Monserrat and Carlos-Jesús Moreno-Ávila.
}

\AbstractEntry
{Tomasz Pełka}
{Fixed point Floer homology: an introduction}
{Jagiellonian University}
{
The course will be a gentle introduction to fixed point Floer theory, aimed at non-experts,
with special focus on possible applications to singularity theory.

First, I will explain Floer's original motivation: the Arnold conjecture, bounding the number
of fixed points of a Hamiltonian diffeomorphism by the sum of Betti numbers. With this in
mind, I will introduce the Floer complex, following the lectures of D. Salamon. Next, I will
explain M. McLean's proposal to use Floer-theoretic invariants to study the geometry of
isolated singularities. Following his ICM address, I will explain why smoothness of isolated
threefold singularities is a symplectic invariant, even though it is not a topological one.

For isolated hypersurface singularities, I will state the arc-Floer conjecture, relating fixed
point Floer homology of the monodromy with the compactly-supported cohomology of contact
loci. This conjecture is motivated by the existence of spectral sequences which converge to
each side and have the same first page. On the Floer side, it is a Morse--Bott type spectral
sequence, due to McLean. I will explain how to construct it from the symplectic version of the
A'Campo model (joint work by J. F. de Bobadilla). Finally, I will indicate some special cases
when this spectral sequence degenerates. In particular, an observation of P. Seidel yields
such degeneration for plane curves, which is one of the ingredients of the proof of arc-Floer
conjectures in this case, due to J. de la Bodega and E. de Lorenzo Poza.
}

\AbstractEntry
{María Martín Vega}
{The Bruno ideal for logarithmic vector fields}
{Institut de Mathématiques de Jussieu--Paris}
{
Given a germ of analytic vector field $\partial$, let
$\partial=\partial_{\rm ss}+\partial_{\rm nilp}$ be its unique formal Jordan decomposition,
where $\partial_{\rm ss}$ is semi-simple and $\partial_{\rm nilp}$ is nilpotent. When
$\partial$ is logarithmic, we define the Bruno ideal $B(\partial)$ as the formal locus where
the semi-simple and the nilpotent components are collinear, that is, the vanishing locus of
$\partial_{\rm ss}\wedge \partial_{\rm nilp}$.

We prove the following result originally stated by A. D. Bruno: if the eigenvalues of
$\partial_{\rm ss}$ satisfy an arithmetic condition, then $B(\partial)$ is analytic and the
restriction of $\partial$ to its zero set $V$ is analytically normalizable.

This talk is based on joint work with Daniel Panazzolo.
}

\AbstractEntry
{Daniel Bath}
{The Strong Monodromy Conjecture and certain homogeneous polynomials}
{KU Leuven}
{
The Strong Monodromy Conjecture promises a miraculous connection between the theory of
algebraic differential operators and the theory of motivic integration, namely it claims the
poles of the motivic zeta function are contained in the zeroes of the Bernstein--Sato
polynomials. Alas, the miracle has few witnesses; even for homogeneous polynomials in three
variables this conjecture is open. We prove it for homogeneous polynomials in three variables
whose associated reduced projective divisor has, at worst, quasi-homogeneous singularities.
In the analogous situation with more variables, we will ``solve'' the D-module part of the
story.

This is joint work with Wim Veys; arXiv:2602.20922.
}

\AbstractEntry
{Mario Morán Cañón}
{Isosingularity in algebraic varieties and arc schemes}
{Universidad Autónoma de Madrid (UAM) / ICMAT}
{
Given a scheme, analytic isosingularity is, informally speaking, the property that the formal
neighborhood of a point remains constant as the point varies. We will present several results,
as well as a number of questions, concerning analytic isosingularity in algebraic varieties on
the one hand, and in arc schemes, that is, the space of formal germs of curves on an algebraic
variety, on the other. This is joint work with David Bourqui.
}

\AbstractEntry
{Roger Gómez López}
{The Archimedean zeta function}
{Universitat Politècnica de Catalunya (UPC)}
{
The Archimedean zeta function is a one-parameter family of distributions, for which the
dependence on the parameter is analytic in the region of convergence. Two approaches have
been used to obtain its meromorphic extension to the complex plane: resolution of
singularities and the Bernstein--Sato polynomial. In the case of irreducible plane curves,
we have obtained a full explicit description of the poles and the residues, and their
consequences for the Bernstein--Sato polynomial.
}

\AbstractEntry
{Alex Hof}
{Hilbert schemes of points and the lattice cohomology of finite maps}
{Alfréd Rényi Institute of Mathematics}
{
The analytic lattice cohomology of a reduced curve germ categorifies its delta invariant and
captures information about other invariants, including its Cohen--Macaulay type and, in the
planar case, its Heegaard Floer link homology. We show that these cohomology groups can be
computed from filtrations on an appropriate collection of nested punctual Hilbert schemes;
this allows us to define a theory of lattice cohomology for finite maps between
arbitrary-dimensional complex-analytic multigerms such that the lattice cohomology of a curve
germ is retrieved as the cohomology of its normalization map. As an application, we construct
specialization maps in lattice cohomology arising from deformations of curve germs. Based on
joint work with András Némethi.
}

\AbstractEntry
{Manuel González Villa}
{Denef--Loeser zeta functions of suspensions and Lê--Yomdin singularities}
{Universidad de Zaragoza}
{
In this talk, we provide formulas for the motivic and topological zeta functions for a broad
family of hypersurfaces. This family includes suspensions of hypersurfaces by an arbitrary
number of points and extends beyond the classical Thom--Sebastiani type. Furthermore, we
present applications to the monodromy and holomorphy conjectures for suspensions and
Lê--Yomdin singularities. This is joint work with Enrique Artal Bartolo, Pedro D. González
Pérez, and Edwin León Cardenal (arXiv:2604.24523v1).
}

\AbstractEntry
{Zsolt Baja}
{Cyclic quotient surface singularities and proportionally modular numerical semigroups}
{Babeș-Bolyai University}
{
Let $(X,o)$ be a weighted homogeneous surface singularity whose link is a rational homology
sphere. By the Pinkham--Dolgachev--Demazure formula, the dimensions of the graded components
of its coordinate ring
\[
R_{(X,o)}=\bigoplus_{\ell\geq 0} R_{(X,o),\ell}
\]
are determined by the topology of the singularity. Therefore one can associate with $(X,o)$
the numerical semigroup
\[
\mathcal{S}
=
\left\{
\ell\in\mathbb{Z}_{\geq 0}
\mid
\dim R_{(X,o),\ell}\neq 0
\right\}.
\]
Numerical semigroups obtained in this way are called \emph{representable numerical
semigroups}. This construction makes it possible to apply tools from singularity theory to the
study of numerical semigroups.

In 2020, László and Némethi derived a formula for the Frobenius number of a representable
numerical semigroup and established the representability of a class of numerical semigroups.
Subsequent works provided a characterization of representable semigroups and established a
formula for their genus.

In this talk, we focus on the case in which $(X,o)$ is a cyclic quotient singularity. The
associated semigroups form the class of \emph{proportionally modular numerical semigroups}
introduced by Rosales et al. We describe this correspondence and illustrate how to determine a
system of generators, as well as how to construct a resolution graph representing a given
semigroup. Finally, we present a solution to a previously open problem posed by Rosales and
García-Sánchez in their 2009 monograph regarding modular numerical semigroups.

This is joint work with T. László and A. Némethi.
}

\AbstractEntry
{Irma Pallarés Torres}
{On the Hodge index theorem for singular spaces and characteristic classes}
{Universidad de Cantabria}
{
The classical Hodge index theorem identifies the signature of a compact complex manifold with
the value at $y=1$ of its Hirzebruch $\chi_y$-genus. Recently, this result has been extended
to singular spaces. Moreover, Brasselet, Schürmann, and Yokura conjectured a characteristic
class version of this theorem. In this talk, we will discuss progress on this conjecture.
This presentation is based on joint work with Fernández de Bobadilla and Saito.
}

\AbstractEntry
{Eduardo de Lorenzo Poza}
{Geometry of contact loci via semihomogeneous singularities}
{BCAM (Basque Center for Applied Mathematics)}
{
Given a scheme $X$, a regular function $f$ and an integer $m$, the $m$-contact locus
$\mathcal{X}_m(f)$ is the subspace of the $m$-jet space $J_m(X)$ consisting of those arcs that
meet the hypersurface $f^{-1}(0)$ with order precisely $m$. Even though the additive
invariants of $\mathcal{X}_m(f)$ are relatively well understood thanks to the tools of
motivic integration, other geometric invariants such as the number of connected or
irreducible components, codimension, cohomology, etc. seem more challenging to comprehend.
In this talk we will explain how these geometric invariants are related to other areas such
as birational geometry and Floer theory, and we will use the example of semihomogeneous
singularities to showcase some unexpected properties of contact loci.
}

\clearpage
\section*{Posters and Lightning Talks}
\addcontentsline{toc}{section}{Posters and Lightning Talks}

\subsection*{Tuesday, June 30}

\begin{itemize}[leftmargin=*]
    \item \emph{Zariski site as a precohesive topos} --- Lidia Angélica Ávila Déciga.
    \item \emph{Real algebraic realization of divide links} --- Arthur Garcia Tonus.
    \item \emph{Holomorphic mappings and their frontalisations} --- Christian Muñoz Cabello.
    \item \emph{Jets of good real images} --- Ignacio Breva Ribes.
    \item \emph{The Singular Derived Category of a Hypersurface} --- Jesús Javier Vidales Pérez.
    \item \emph{Essential divisors and the higher Nash modification for RDP singularities} --- José Alejandro Tenorio Vázquez.
\end{itemize}

\subsection*{Thursday, July 2}

\begin{itemize}[leftmargin=*]
    \item \emph{Higher-order connections for modules} --- Jesús Daniel Aparicio Barrera.
    \item \emph{Cubical hyperresolutions for stable corank $\leq 1$ maps} --- José Galindo Jiménez.
    \item \emph{Line arrangements on Cubic Surfaces} --- Juan Carlos Castro Rivera.
    \item \emph{The Link of Polar Weighted Homogeneous Polynomials} --- Jessica Torres Flores.
    \item \emph{On the isotropic points of holomorphic curves} --- Amanda Dias Falqueto.
\end{itemize}

\end{document}
